\documentclass{article}
\usepackage[utf8]{inputenc}
\usepackage[T1]{fontenc}
\usepackage{geometry}
\geometry{
    a4paper,
    left=2.5cm,
    right=2.5cm,
    top=2.5cm,
    bottom=2.5cm
}
\usepackage{graphicx}
\usepackage{float}
\usepackage{hyperref}
\usepackage[english]{babel}
\usepackage{amsmath}
\usepackage{amssymb}
\usepackage{array}
\usepackage{booktabs}
\usepackage{caption}
\usepackage{subcaption}
\usepackage{enumitem}
\usepackage{fancyhdr}
\usepackage{titlesec}
\usepackage{url}
\usepackage{textcomp}
\usepackage{xcolor}
\usepackage{listings}
\usepackage{multirow}
\usepackage{longtable}
\usepackage{tabularx}
\usepackage{times}

% Set up page style
\pagestyle{fancy}
\fancyhf{}
\rfoot{\thepage}
\renewcommand{\headrulewidth}{0pt}

% Title formatting
\titleformat{\chapter}[display]{\normalfont\bfseries}{\chaptertitle\thechapter}{0pt}{\Huge}
\titlespacing*{\chapter}{0pt}{50pt}{40pt}

% Document info
\title{Template Document}
\author{Academic Support Agent}
\date{\today}

\begin{document}
\maketitle

%% preamble page numbers with roman numerals
\pagenumbering{roman}\setcounter{page}{1}
\pagestyle{plain}

%%------------------------------- cover page ----------------------------------

\begin{titlepage}
\center

\vspace{-15mm}
{\large \textbf{\textsc{\school}}}\\

\vfill

{\Large \textbf{\projtitle}}\\[8mm]
{\large \textbf{\subtitle}}\\[28mm]

{\Large \textbf{\projauthor}}\\

\vfill

\includegraphics[width=52mm]{uporto-feup.pdf}

\vfill

{\large \degree}\\[8mm]
{\large \textbf{Tutor na U.Porto}: \tutor}\\[2mm]
{\large \textbf{Orientador na empresa}: \supervisor}\\[8mm]

%\renewcommand{\today}{15 de dezembro de 2023}
\today

\end{titlepage}

%%------------------------------- Abstract ------------------------------------

\chapter*{Resumo}

\begin{info}
O resumo tem um car ter essencialmente informativo, devendo ser
escrito de forma concisa (at  200 palavras) de maneira a captar o
interesse de quem o vai ler.

O Resumo substitui a leitura do documento e n o contem figuras,
tabelas, cita es, etc.\ 
Deve incluir os seguintes t picos:  mbito, objetivos, os m todos, as
principais descobertas, incluindo resultados, conclus es e
recomenda es, se existirem.

Para saber mais sobre como redigir um bom resumo consulte o tutorial
online dispon vel no website na Biblioteca, ``Guia de Apoio  
Publica o'', sec o: 
``\href{https://docs.google.com/document/d/1TDC1behVq8x7fQL4CcPEEh_np5GXviJevQxnQ9gbiJs/edit\#heading=h.s4z9k57ywd9w}
{Estruturar Relat rio T cnico}''.
\end{info}

\todo[inline]{Escrever o Resumo, mas s  no fim.}

%\vspace{\fill}
%{\Large \textbf{Palavras-chave}:} palavra1, palavra2, palavra3, palavra4
%\vspace*{24mm}

%%------------------------------- Acknowledgments -----------------------------

\chapter*{Agradecimentos}

\begin{infoopt}
Habitualmente   mencionada a contribui o de outras pessoas ou
entidades, tanto para a realiza o do estudo como para a produ o do
relat rio. 
Podem fazer-se numa p gina aut noma ou incluir-se na introdu o.
\end{infoopt}

%%------------------------------- table of contents ---------------------------

%% redefine tableofcontents text, ONLY for PT
\renewcommand{\contentsname}{ ndice}



%%------------------------------- list of todos -------------------------------

% list todos; comment in the end (should be empty before delivery :-)
\listoftodos

\begin{infoopt}
Podem ser colocadas anota es durante a prepara o do documento, que
s o listadas aqui.

Este elemento n o aparece no documento final!
\end{infoopt}

%%------------------------------- Glossary ------------------------------------

\chapter*{Gloss rio}
%\addcontentsline{toc}{chapter}{Gloss rio}

\begin{description}
\item[bash] \hfill \\
  Bash   uma \emph{shell Unix} e uma linguagem de comando escrita 
  em 1989 por Brian Fox para o Projeto GNU como um substituto de 
  software livre para a \emph{Bourne shell}.
\item[firewall] \hfill \\
  Em computa o, uma \emph{firewall}   um sistema de seguran a de rede 
  que monitoriza e controla o tr fego de entrada e sa da da rede 
  com base em regras de seguran a predeterminadas.
  Uma \emph{firewall} normalmente estabelece uma barreira entre uma 
  rede confi vel e uma rede n o confi vel, como a Internet.
\item[Gloss rio] \hfill \\
  Gloss rio   uma esp cie de pequeno dicion rio espec fico para
  palavras e express es pouco conhecidas presentes num texto, seja
  por serem de natureza t cnica, regional ou de outro idioma.
\end{description}

\begin{infoopt}
Justifica-se ter um gloss rio sempre que seja necess rio esclarecer o leitor sobre o significado de terminologia espec fica usada no texto no relat rio.
Recomenda-se a sua localiza o nos elementos iniciais, embora na
normaliza o existente haja variantes, podendo tamb m constar nos
elementos finais.

A lista de itens deve ser ordenada alfabeticamente\footnote{%
Para saber mais consulte o tutorial online 
``\href{https://docs.google.com/document/d/1TDC1behVq8x7fQL4CcPEEh_np5GXviJevQxnQ9gbiJs/edit}
{Guia de Apoio   Publica o}''.}.
\end{infoopt}

%%------------------------------- chapter ------------------------------------

\chapter{Introdu o}

%% display headers & footers
\pagestyle{fancy}
%% main page numbers with arabic numerals
\pagenumbering{arabic}\setcounter{page}{1}

\begin{info}
Contextualiza o sucinta do assunto do relat rio, fazendo-se
refer ncia ao  mbito e aos objetivos.
Aqui clarifica-se a motiva o do trabalho apresentado e explica-se a
abordagem adotada e a sua rela o com trabalhos an logos, numa
perspetiva gen rica.
N o se deve antecipar detalhes sobre o que   explicado nas sec es
posteriores. 
Se for pertinente, pode-se indicar ainda qual o p blico a que se
destina\footnote{%
Para saber mais consulte o tutorial online 
``\href{https://docs.google.com/document/d/1TDC1behVq8x7fQL4CcPEEh_np5GXviJevQxnQ9gbiJs/edit}
{Guia de Apoio   Publica o}''.}.
\end{info}

Este relat rio aborda \ldots\ fornece uma vis o geral do contexto, objetivos e abordagem adotada.

\section{Contexto}

\begin{info}
   Apresentar o contexto organizacional em que decorreu o projeto/est gio (empresa, institui o, unidade de investiga o, laborat rio de investiga o, etc.). 
   Apresentar o problema abordado e a motiva o para o trabalho realizado (qual   o problema abordado e porque   que   importante).
\end{info}

\section{Objetivos e resultados esperados}

\begin{info}
  Indicar os objetivos do trabalho e os resultados esperados.
\end{info}

\section{Estrutura do relat rio}
	
\begin{info}
  Descrever brevemente a estrutura do relat rio.

  Ser  expect vel que o relat rio tenha entre 10 a 25 p ginas (em formato A4, coluna  nica, com um tamanho de letra n o excedendo os 12pt no texto de par grafo), j  contando com anexos. 
\end{info}

Para al m da introdu o, este relat rio est  organizado em mais 4 cap tulos.
No Cap tulo~\ref{chap:metodo} \ldots

%%------------------------------- DELETEME ----------------------------------

\newpage  % don't do this
\section*{Exemplos}

\todo[inline]{Remover a sec o ``Exemplos'', quando j  n o for necess ria.}

\begin{info}
  S o ilustradas de seguida algumas partes do documento.

  Esta sec o n o aparece no documento final!
\end{info}

\subsection*{Equa es}

\begin{info}
Este texto   apenas um exemplo que precede uma equa o.
\end{info}  

Equa es simples podem ser inseridas em linha com o texto: 
a reta   \(y=mx+b\).

Equa es mais complicadas devem ser separadas em linhas individuais e
numeradas sequencialmente   direita dentro de par ntesis.
Esta   a equa o 2  grau gen rica:

\begin{equation} \label{eq:1}
  ax^2+bx+cx=0
\end{equation}

Onde $a$   o coeficiente de 2  grau; $b$ o de 1  grau; $c$ o
coeficiente independente da vari vel $x$, a determinar.

As equa es devem ser referidas mantendo o seu n mero.
Por exemplo, a Equa o~\ref{eq:2} resolve problemas formulados tal como 
mostrado na Equa o~\ref{eq:1}.

\begin{equation} \label{eq:2}
  x=\frac{-b\pm \sqrt{b^2-4ac}}{2a}
\end{equation}

\subsection*{Figuras e tabelas}

Todas as figuras e tabelas devem ser obrigatoriamente legendadas e
numeradas sequencialmente:

\begin{itemize}
\item as figuras devem ser legendadas por baixo;
\item as tabelas devem ser legendadas no topo. 
\end{itemize}

Mantenha as figuras centradas e em linha com o texto para que a
legenda apare a sempre colada com a imagem.

\begin{info}
As figuras devem flutuar livremente na p gina e ser referidas e
descritas no texto, com as fontes devidamente explicitadas, para
evitar o pl gio.
\end{info}

Como exemplo, a Figura~\ref{fig:campus} % (retirada de \url{www.fe.up.pt}) 
mostra o \emph{campus} da FEUP. 

\lipsum[1]

\begin{figure}
\centering
\includegraphics[width=0.8\textwidth]{campus}
\caption[Fotografia a rea do Campus da FEUP]{Fotografia a rea do Campus da FEUP~\cite{kn:figura}.} \label{fig:campus}
\end{figure}

\lipsum[2]

\begin{info}
Pode ser reservado espa o para colocar, no futuro, uma figura; por
exemplo, a Figura~\ref{fig:natal}.
\end{info}

\lipsum[3]

\begin{figure} [b]
  \centering
  \missingfigure{Inserir a figura do Natal.}
  \caption{O Natal no Campus da FEUP.} \label{fig:natal}
\end{figure}

\lipsum[4]

\begin{info}
As tabelas devem flutuar livremente na p gina e ser referidas e
descritas no texto, com as fontes devidamente explicitadas, para
evitar o pl gio.
\end{info}

A Tabela~\ref{tab:feup} % (excerto adaptado de ``A FEUP em n meros'', 2011)
serve para exemplificar como mostrar alguns valores que, neste caso, t m
a ver com alguns dados num ricos associados a recursos e investimentos
da FEUP no ano de 2011. 

\lipsum[4]

\begin{table}
  \centering
  \caption[Recursos F sicos da FEUP]{Recursos F sicos da FEUP. Adaptado de~\cite{kn:tabela}}
  \begin{tabular}{| l | c |}
    \hline
    \textbf{Designa o} & \textbf{Quantidade}\\\hline
    \hline
     rea total do campus FEUP & 93 918 $m^2$\\\hline
    Espa os verdes & 23 000 $m^2$\\\hline
    N mero de computadores dedicados ao ensino & 1 815\\\hline
    Investimento em equipamentos de laborat rio & 1,46 M \\
    \hline
  \end{tabular}
  \label{tab:feup}
\end{table}

\lipsum[6]

\subsection*{Cita es}

  medida que escreve o texto do relat rio, deve indicar os trabalhos
de outros autores em que se baseia, sob a forma de cita es.
Isto consiste em indicar de forma abreviada as fontes usadas  s quais
foi buscar informa o adicional para desenvolver o tema do seu
relat rio.

Existem duas formas principais de citar:
\begin{itemize}
\item por \textbf{par frase}: interpreta o do conte do original por
  palavras diferentes das da fonte consultada, indicando a fonte logo
  a seguir; ou 
\item 
  por \textbf{transcri o}: uso de um excerto do conte do original
  apresentando-o entre aspas, indicando a fonte logo a seguir.
\end{itemize}

As cita es devem obedecer a um estilo normalizado.
De entre os muitos que existem, a Biblioteca da FEUP aconselha o
estilo Chicago (formato autor-data).

\begin{info}
De seguida exemplificam-se, ao acaso, algumas cita es (por par frase)
de acordo com esse estilo.
\end{info}

A decis o de escolha de um tema para um trabalho acad mico pode
variar~\cite{kn:Bel02-book}.
O tema pode ser pensado e escolhido pelo pr prio estudante, ou a
partir de uma lista de temas j  concebidos, com potencial interesse
para estudo~\cite{kn:GLPR14-joPhysics}.

A cada cita o ao longo do texto deve corresponder uma refer ncia
indicada na lista final de refer ncias
bibliogr ficas~\cite{kn:Lip08,kn:MSS+12-wemep,kn:VKL+18-dtu}. 

  importante n o esquecer que tamb m as figuras (imagens, tabelas,
gr ficos, etc.) provenientes de obras de outros autores (por exemplo 
obtidas atrav s da Internet) devem ser citadas sempre, ap s as
respetivas legendas~\cite{kn:GLPC22-torque}.

Para saber mais sobre este assunto e ver exemplos, consulte o guia
``Evitar o pl gio: boas pr ticas no uso da
informa o''\footnote{\url{https://feup.libguides.com/plagio/citar}}.  

Duis non odio morbi ut dui sed accumsan risus eget odio~\cite{iso_19156_2011,ornelas2016}. 

%\lipsum[7]

\subsection*{C digo}

\begin{info}
De seguida   ilustrada a inclus o de c digo no documento.
\end{info}

Listing~\ref{code:useless} 
\lipsum[8]

\begin{lstlisting}[language=Python, caption=Python example, label=code:useless]
# Take the user's input
words = input("Enter the text to translate to pig latin: ")
print(f"You entered: {words}")

# Break apart the words into a list
words = words.split(' ')

# Use a list comprehension to translate words greater than or equal to 3 characters
translated_words = [(w[1:] + w[0] + "ay") for w in words if len(w) >= 3 ]

# Print each translated word
for word in translated_words:
    print(word)
\end{lstlisting}

\subsection*{Uso das macros}

\begin{info}
De seguida   ilustrada a utliza o de macros \LaTeX{} definidas no
pre mbulo.
De salientar o uso da macro \verb!\class{}! para classes, m todos e componentes.
\end{info}

O \windspt, retirado do \windscannerpt, usa o \svg\ \ldots\ como em \class{Student.calculate-age()}.

\lipsum[9]

\begin{info}
As partes componentes subsequentes que constituem o corpo do texto
devem ser estruturadas em sec es, estimando-se que at  3 n veis seja
o suficiente para este tipo de trabalho.

Para saber mais consulte o tutorial online 
``\href{https://docs.google.com/document/d/1TDC1behVq8x7fQL4CcPEEh_np5GXviJevQxnQ9gbiJs/edit}
{Guia de Apoio   Publica o}''.
Note-se que as se es a  indicadas podem ser adaptadas em fun o do tema
ou profundidade do estudo a desenvolver.
\end{info}

\begin{info}
N o   costume haver cabe alhos de sec es seguidas sem texto.
\end{info}

\subsection*{O travess o}

Sobre o uso do h fen e do travess o\footnote{"Ciberd vidas da L ngua Portuguesa", \url{https://ciberduvidas.iscte-iul.pt/consultorio/perguntas/o-uso-do-hifen-e-do-travessao/31251}}: 
\begin{enumerate}
    \item O h fen (Alt + 0045): sem espa os em branco (-);
    \item O travess o (ou travess o duplo)/em-dash(1) (Alt + 0151): em portugu s, envolto em espa os em branco (---);
    \item O tra o de liga o/en dash (Alt + 150): sem espa os em branco (--);
    \item O sinal matem tico de subtra o: sem espa os em branco (--).
\end{enumerate}

\subsection*{As aspas}

Sobre as aspas  em \LaTeX, ou   usado o slifo diretamente, ou ent o s o feitas com o acento para tr s no in cio e a plica no fim, como em ``exemplo''\footnote{What is the best way to use quotation mark glyphs? 
\url{https://tex.stackexchange.com/questions/531/what-is-the-best-way-to-use-quotation-mark-glyphs}}.

%%------------------------------- chapter ------------------------------------

\chapter{Metodologia utilizada e principais atividades desenvolvidas}\label{chap:metodo}

Neste cap tulo   descrita a metodologia seguida, s o enumerados os
principais intervenientes no projeto e   feito o registo das
principais atividades desenvolvidas.

\section{Metodologia utilizada}

\begin{info}
Descrever a metodologia~\cite{despa2014comparative} seguida (por
exemplo, desenvolvimento iterativo com \emph{sprints} quinzenais e reuni es
semanais de acompanhamento) e os recursos utilizados (por exemplo,
GitHub\footnote{\url{https://github.com/}}, etc.). 
\end{info}

\section{Intervenientes, pap is e responsabilidades}

\begin{info}
Identificar a equipa de projeto, \emph{stakeholders} e outros invernientes com os quais existiu intera o; no caso de trabalho em grupo, clarificar os pap is e responsabilidades de cada elemento do grupo.
\end{info}

\section{Atividades desenvolvidas}

\begin{info}
Descrever as atividades realizadas ao londo do tempo (incluindo eventos relevantes, como apresenta es, reuni es com clientes, etc.) e respectivos \emph{deliverables}, recorrendo tipicamente a um diagrama de Gantt~\cite{gantt} (ver Figura~\ref{fig:gantt}) e a uma descri o sum ria de cada atividade/\emph{deliverable}.
Pode ser apresentado tamb m atrav s de tabela com progresso semanal.
\end{info}

\begin{figure}
    \begin{ganttchart}[vgrid, hgrid, title height=1,
        x unit=6mm, y unit title=8mm, y unit chart=7mm]{1}{21}

        \gantttitle{Fev.}{4}
        \gantttitle{Mar.}{4}
        \gantttitle{Abr.}{5}
        \gantttitle{Mai.}{4}
        \gantttitle{Jun.}{4} \\

        \gantttitle{}{1}
        \gantttitle{S1}{1}
        \gantttitle{S2}{1}
        \gantttitle{S3}{1}
        \gantttitle{S4}{1}
        \gantttitle{S5}{1}
        \gantttitle{S6}{1}
        \gantttitle{S7}{1}
        \gantttitle{P}{1}
        \gantttitle{S8}{1}
        \gantttitle{S9}{1}
        \gantttitle{S10}{1}
        \gantttitle{S11}{1}
        \gantttitle{Q}{1}
        \gantttitle{S12}{1}
        \gantttitle{S13}{1}
        \gantttitle{S14}{1}
        \gantttitle{S15}{1}
        \gantttitle{S16}{1}
        \gantttitle{S17}{1}
        \gantttitle{S18}{1} \\

        \ganttbar{\textbf{Reuni o kick-off}}{1}{1} \\
        \ganttgroup{Trabalho na empresa}{2}{17} \\
        \ganttbar{\textbf{Reuni o interm dia}}{11}{11} \\
        \ganttbar{\textbf{Reuni o final}}{18}{18} \\
        \ganttbar{Escrita do Relat rio}{18}{20} \\
        \ganttbar{Entregas}{21}{21}
    \end{ganttchart}
    \caption{Diagrama de Gantt.} \label{fig:gantt}
\end{figure}

\chapter{Desenvolvimento da solu o}

Neste cap tulo   descrito o trabalho desenvolvido para alcan ar os
resultados esperados. 

Se for o caso de um prot tipo de software, s o apresentados os
requisitos, a arquitetura da solu o, o desenvolvimento e a valida o
do prot tipo.

\section{Requisitos}

\begin{info}
Identificar os requisitos funcionais e n o funcionais relevantes e respectivas fontes, bem como restri es ao projeto.

Se seguida apresenta-se um exemplo de tabela com \textit{user stories}, que pode ser usada no caso do n mero de linhas ser grande.
Pode ser necess ria mais uma coluna para prioridades.
\end{info}

\subsection*{User stories}

Exemplos de poss veis cen rios de utiliza o s o apresentados na Tabela~\ref{tab:userstories}, utilizando \textit{user stories} escritas de acordo com o modelo: 
``\textbf{As a <type of user>, I want <some goal> so that <some reason>}''.

\begin{center}
\begin{longtable}{| c | p{18mm} | p{110mm} |}
\caption{Usage scenario} \label{tab:userstories} \\[-2mm] \hline

\textbf{Identifier} & \textbf{Name} & \textbf{Description} \\ \hline\hline
\endfirsthead

\multicolumn{3}{c}{{\tablename\ \thetable{} -- continued from previous page}} \\[2mm] \hline
\textbf{Identifier} & \textbf{Name} & \textbf{Description} \\ \hline\hline
\endhead

\multicolumn{3}{r}{{Continues on next page\ldots}} \\
\endfoot

\hline
\endlastfoot

US01 & Register (High) & As a \textit{Visitor}, I want to register a new account so that I have access to profile data  \\ \hline

% US02 & Login (High) & As a \textit{Visitor}, I want to login so that % I have access to profile data  \\ \hline
% 
% US03 & Edit profile (High) & As a \textit{Authenticated user}, I want % to edit \ldots\  \\ \hline

\end{longtable}
\end{center}

\section{Arquitetura e tecnologias (ou Conce o e Implementa o)}

\begin{info}
Arquitetura e tecnologias utilizadas e respetiva justifica o (tendo em conta requisitos e alternativas existentes), diagramas t cnicos elaborados (ver Figura~\ref{fig:componentes} retirada de SPARX\footnote{\url{https://sparxsystems.com/resources/tutorials/uml2/component-diagram.html}}), dificuldades t cnicas encontradas e sua resolu o, etc.
\end{info}

\begin{figure}
\centering
\includegraphics[width=0.72\textwidth]{components}
\caption{Diagrama de Componentes em UML.} \label{fig:componentes}
\end{figure}

\section{Solu o desenvolvida}

\begin{info}
Apresentar a solu o desenvolvida na  tica do utilizador, com ajuda de \emph{screenshots}.
\end{info}

\section{Valida o}

\begin{info}
Descri o da valida o da solu o desenvolvida, em rela o aos requisitos e restri es identificados, e respectivos resultados (por exemplo, resultados de avalia o experimental, testes efetuados, \emph{feedback} de utilizadores ou especialistas, etc.).
\end{info}

\chapter{Conclus es}

Neste cap tulo s o sumariados os resultados alcan ados e as li es
aprendidas.
Finalmente, s o apresentadas as limita es do trabalho e s o propostas
melhorias e trabalho futuro.

\section{Resultados alcan ados}

\begin{info}
Sumariar os resultados alcan ados e contribui es (em rela o aos  objetivos).

No caso de trabalho em grupo, clarificar as contribui es individuais, em termos qualitativos e quantitativos (percentagem).
\end{info}

\lipsum[11]

\section{Li es aprendidas}

\begin{info}
Refletir sobre as li es aprendidas (tendo em conta os objetivos de aprendizagem).
\end{info}

\lipsum[12]

\section{Trabalho futuro}

\begin{info}
Identificar limita es do trabalho realizado e ideias de melhorias e trabalho futuro.
\end{info}


%%------------------------------- Bibliography --------------------------------

%Aqui dever o constar as refer ncias bibliogr ficas constantes do texto do relat rio e eventual bibliografia relevante consultada durante o projeto.

\renewcommand{\bibname}{Refer ncias bibliogr ficas}
\printbibliography
\addcontentsline{toc}{chapter}{\refname}  % add it to table of contents

\begin{info}
Na lista final de refer ncias devem constar os trabalhos dos autores
citados de forma abreviada ao longo do texto, obtida automaticamente
com o \class{BibTeX}.
A refer ncia bibliogr fica   a forma mais desenvolvida de indicar as
fontes de informa o em que se baseou. 
\end{info}

%%------------------------------- appendix ------------------------------------


\chapter{Um Ap ndice}

\begin{infoopt}
Os ap ndices e os anexos cont m informa o que complementa, apoia e
clarifica o relat rio e cuja inclus o no corpo principal do relat rio
interferiria com uma boa ordem de apresenta o das ideias.

H  uma diferen a importante entre ap ndices e anexos: ``No ap ndice
compilam-se apenas os documentos que s o da autoria do autor do
relat rio, enquanto no anexo se compilam os documentos de autoria de
outros autores que n o o do relat rio.''\footnote{%
Para saber mais consulte o tutorial online 
``\href{https://docs.google.com/document/d/1TDC1behVq8x7fQL4CcPEEh_np5GXviJevQxnQ9gbiJs/edit}
{Guia de Apoio   Publica o}''.}
\end{infoopt}

%%------------------------------- the end. ------------------------------------
\end{document}